\SourceLeaf{248}%
pacto, si variatio per characteristicam \(\delta\) denotatur, habemus
\[
\begin{aligned}
\delta s&=\int\frac{(\tfrac{\mathrm{d}E}{\mathrm{d}p}\mathord{.}\mathrm{d}p^2+\tfrac{2\mathrm{d}F}{\mathrm{d}p}\mathord{.}\mathrm{d}p\mathord{.}\mathrm{d}q+\tfrac{\mathrm{d}G}{\mathrm{d}p}\mathrm{d}q^2)\delta p+(2E\mathrm{d}p+2F\mathrm{d}q)\mathrm{d}\delta p}{2\mathrm{d}s}\\
&=\frac{E\mathrm{d}p+F\mathrm{d}q}{\mathrm{d}s}\mathord{.}\delta p+\int\delta p\mathord{.}\{\frac{\tfrac{\mathrm{d}E}{\mathrm{d}p}\mathord{.}\mathrm{d}p^2+\tfrac{2\mathrm{d}F}{\mathrm{d}p}\mathord{.}\mathrm{d}p\mathord{.}\mathrm{d}q+\tfrac{\mathrm{d}G}{\mathrm{d}p}\mathord{.}\mathrm{d}q^2}{2\mathrm{d}s}-\mathrm{d}\mathord{.}\frac{E\mathrm{d}p+F\mathrm{d}q}{\mathrm{d}s}\}
\end{aligned}
\]
constatque, quae hic sunt sub signo integrali, independenter a \(\delta p\) evanescere debere. Fit itaque
\[
\begin{aligned}
\frac{\mathrm{d}E}{\mathrm{d}p}\mathord{.}\mathrm{d}p^2+\frac{2\mathrm{d}F}{\mathrm{d}p}\mathord{.}\mathrm{d}p\mathord{.}\mathrm{d}q+\frac{\mathrm{d}G}{\mathrm{d}p}\mathord{.}\mathrm{d}q^2&=2\mathrm{d}s\mathord{.}\mathrm{d}\mathord{.}\frac{E\mathrm{d}p+F\mathrm{d}q}{\mathrm{d}s}\\
&=2\mathrm{d}s\mathord{.}\mathrm{d}\mathord{.}\BareRoot{E}\mathord{.}\cos\theta\\
&=\frac{\mathrm{d}s\mathord{.}\mathrm{d}E\mathord{.}\cos\theta}{\BareRoot{E}}-2\mathrm{d}s\mathord{.}\mathrm{d}\theta\mathord{.}\BareRoot{E}\mathord{.}\sin\theta\\
&=\frac{(E\mathrm{d}p+F\mathrm{d}q)\mathrm{d}E}{E}-\BareRoot{(EG-FF)}\mathord{.}\mathrm{d}q\mathord{.}\mathrm{d}\theta\\
&=(\frac{E\mathrm{d}p+F\mathrm{d}q}{E})\mathord{.}(\frac{\mathrm{d}E}{\mathrm{d}p}\mathord{.}\mathrm{d}p+\frac{\mathrm{d}E}{\mathrm{d}q}\mathord{.}\mathrm{d}q)\\
&\quad-2\BareRoot{(EG-FF)}\mathord{.}\mathrm{d}q\mathord{.}\mathrm{d}\theta
\end{aligned}
\]
Hinc itaque nanciscimur aequationem conditionalem pro linea brevissima sequentem:
\[
\begin{aligned}
\BareRoot{(EG-FF)}\mathord{.}\mathrm{d}\theta
&=\tfrac{1}{2}\frac{F}{E}\mathord{.}\frac{\mathrm{d}E}{\mathrm{d}p}\mathord{.}\mathrm{d}p+\tfrac{1}{2}\mathord{.}\frac{F}{E}\mathord{.}\frac{\mathrm{d}E}{\mathrm{d}q}\mathord{.}\mathrm{d}q\\
&\quad+\tfrac{1}{2}\mathord{.}\frac{\mathrm{d}E}{\mathrm{d}q}\mathord{.}\mathrm{d}p-\frac{\mathrm{d}F}{\mathrm{d}p}\mathrm{d}p-\tfrac{1}{2}\mathord{.}\frac{\mathrm{d}G}{\mathrm{d}p}\mathrm{d}q
\end{aligned}
\]
quam etiam ita scribere licet
\[
\BareRoot{(EG-FF)}\mathord{.}\mathrm{d}\theta=\tfrac{1}{2}\frac{F}{E}\mathord{.}\mathrm{d}E+\tfrac{1}{2}\mathord{.}\frac{\mathrm{d}E}{\mathrm{d}q}\mathord{.}\mathrm{d}p-\frac{\mathrm{d}F}{\mathrm{d}p}\mathord{.}\mathrm{d}p-\tfrac{1}{2}\mathord{.}\frac{\mathrm{d}G}{\mathrm{d}p}\mathord{.}\mathrm{d}q
\]
Ceterum adiumento aequationis
\[
\mathrm{cotg}\,\theta=\frac{E}{\BareRoot{(EG-FF)}}\mathord{.}\frac{\mathrm{d}p}{\mathrm{d}q}+\frac{F}{\BareRoot{(EG-FF)}}
\]
ex illa aequatione angulus \(\theta\) eliminari, atque sic aequatio differentio-differentialis inter \(p\) et \(q\) evolvi potest, quae tamen magis complicata et ad applicationes minus utilis evaderet, quam praecedens.
\Article{19}
Formulae generales, quas pro mensura curvaturae et pro variatione directionis lineae brevissimae in artt. 11, 18 eruimus, multo simpliciores fiunt, si quantitates \(p,q\) ita sunt electae, ut lineae primi systematis lineas secundi systematis ubique orthogonaliter secent, i. e. ut generaliter habeatur \(\omega=90^0\), sive \(F=0\). Tunc scilicet fit, pro mensura curvaturae,
